% GENERATED FILE. Edit canonical lessons, then run npm run generate:books.
% canonical-source-hash: fnv1a64:f3e900970c4bdb80
% canonical-lessons: TE-C324-vaste, TE-C324-varsham-padite, TE-R324-if-and-otherwise

\chapter{If It Rains}
\label{ch:te-if-it-rains}
\hlchaptermodality{324}

\begin{chapteropening}
\textbf{By the end of this chapter:} \emph{I can say what will happen if something happens, with -tē and okavēḷa, and if it does not, with -akapōtē and lēkapōtē.}

Complete TE-R324-if-and-otherwise: make real conditions with -tē, okavēḷa, -akapōtē and lēkapōtē.
\end{chapteropening}

\section[\texorpdfstring{\te{వస్తే} (vastē)}{వస్తే (vastē)}]{\te{వస్తే} (vastē) — if}

\label{lesson:TE-C324-vaste}



\begin{quote}
Say \textquotedblleft{}I will go\textquotedblright{} (\emph{nēnu veḷtānu}) and \textquotedblleft{}you had better go to the hospital\textquotedblright{} (\emph{mīru āsupatriki veḷtē mañcidi}). Then \textquotedblleft{}in that case\textquotedblright{} and \textquotedblleft{}if by chance\textquotedblright{}. (\emph{ayitē}, \emph{okavēḷa}.)
\end{quote}

\subsection*{You'll want to know: -\te{తే}, 'if'}
\begin{quote}
\textbf{\te{నువ్వు} \te{వస్తే} \te{నేను} \te{కూడా} \te{వస్తాను}.}
\end{quote}

\begin{quote}
\emph{nuvvu vastē nēnu kūḍā vastānu.} — \textquotedblleft{}If you come, I will come too.\textquotedblright{}
\end{quote}

Where the form you know ends in \textbf{-\te{తాను}}, \textbf{-\te{తే}} takes the place of the ending: \textbf{\te{వస్తాను}} $\to$ \textbf{\te{వస్తే}}, \textbf{\te{వెళ్తాను}} $\to$ \textbf{\te{వెళ్తే}}, \textbf{\te{చేస్తాను}} $\to$ \textbf{\te{చేస్తే}}, \textbf{\te{ఉంటాను}} $\to$ \textbf{\te{ఉంటే}}, \textbf{\te{తింటాను}} $\to$ \textbf{\te{తింటే}}. Some verbs put \textbf{-\te{ఇతే}} on the past stem instead: \textbf{\te{చదివితే}}, \textbf{\te{పడితే}}, \textbf{\te{అడిగితే}}.

It does not change for the person: \textbf{\te{నేను} \te{వస్తే}}, \textbf{\te{అతను} \te{వస్తే}}, \textbf{\te{వాళ్ళు} \te{వస్తే}}. The \textbf{-\te{తే} \te{మంచిది}} of advice was this ending. So is \textbf{\te{అయితే}}, \textquotedblleft{}in that case\textquotedblright{}: the \textbf{-\te{తే}} of the verb \textquotedblleft{}to be\textquotedblright{}, \textquotedblleft{}if it is so\textquotedblright{}.

\subsection*{You'll want to know: \te{ఒకవేళ} … -\te{తే}}
\begin{quote}
\textbf{\te{ఒకవేళ} \te{వర్షం} \te{పడితే}, \te{మనం} \te{ఇంట్లోనే} \te{ఉందాం}.}
\end{quote}

\begin{quote}
\emph{okavēḷa varṣaṁ paḍitē, manaṁ iṇṭlōnē undāṁ.} — \textquotedblleft{}If by chance it rains, let's stay at home.\textquotedblright{}
\end{quote}

\textbf{\te{ఒకవేళ}} in front warns that the condition may well not happen; the \textbf{-\te{తే}} still does the work. \textbf{\te{మనం}} (\emph{manaṁ}, \textquotedblleft{}we, you and I\textquotedblright{}) and \textbf{\te{ఉందాం}} (\emph{undāṁ}, \textquotedblleft{}let's stay\textquotedblright{}), with the \textbf{-\te{దాం}} of \textbf{\te{కలుద్దాం}}. \textbf{\te{ఇంట్లోనే}} (\emph{iṇṭlōnē}, \textquotedblleft{}right at home\textquotedblright{}) leans on \textbf{\te{ఇంట్లో}} with an \textbf{-\te{ఏ}}.

\subsection*{Your turn}
Say these aloud:

\begin{itemize}
\raggedright
  \item \textquotedblleft{}If you come, I will come too\textquotedblright{} — \emph{nuvvu vastē nēnu kūḍā vastānu}
  \item \textquotedblleft{}If he eats, I will eat too\textquotedblright{} — \emph{atanu tiṇṭē nēnu kūḍā tiṇṭānu}
  \item \textquotedblleft{}If by chance it rains, let's stay at home\textquotedblright{} — \emph{okavēḷa varṣaṁ paḍitē, manaṁ iṇṭlōnē undāṁ}
  \item If you read this book, you will like it
\end{itemize}

\begin{tcolorbox}[breakable,colback=teal!4,colframe=teal!35!black,title={Before you move on}]
\textquotedblleft{}If you go\textquotedblright{}? (\textbf{\te{నువ్వు} \te{వెళ్తే}}.) \textquotedblleft{}If there is\textquotedblright{}? (\textbf{\te{ఉంటే}}.) \textquotedblleft{}If it rains\textquotedblright{}? (\textbf{\te{వర్షం} \te{పడితే}}.) Which word you knew is already a \textbf{-\te{తే}} form? (\textbf{\te{అయితే}}.)
\end{tcolorbox}
\section[\texorpdfstring{\te{వర్షం} \te{పడితే} (varṣaṁ paḍitē)}{వర్షం పడితే (varṣaṁ paḍitē)}]{\te{వర్షం} \te{పడితే} (varṣaṁ paḍitē) — if it rains, and if not}

\label{lesson:TE-C324-varsham-padite}



\begin{quote}
Say \textquotedblleft{}if you come, I will come too\textquotedblright{} (\emph{nuvvu vastē nēnu kūḍā vastānu}). Then \textquotedblleft{}otherwise\textquotedblright{}. (\emph{lēkapōtē}.)
\end{quote}

\subsection*{You'll want to know: a real condition}
\begin{quote}
\textbf{\te{రేపు} \te{వర్షం} \te{పడితే} \te{నేను} \te{రాను}.}
\end{quote}

\begin{quote}
\emph{rēpu varṣaṁ paḍitē nēnu rānu.} — \textquotedblleft{}If it rains tomorrow, I will not come.\textquotedblright{}
\end{quote}

\begin{quote}
\textbf{\te{సమయం} \te{ఉంటే} \te{సినిమాకి} \te{వెళ్దాం}.}
\end{quote}

\begin{quote}
\emph{samayaṁ uṇṭē sinimāki veḷdāṁ.} — \textquotedblleft{}If there is time, let's go to a film.\textquotedblright{}
\end{quote}

A real condition is one that may well happen. The main verb is the form you know for what will happen (\textbf{\te{వస్తాను}}, \textbf{\te{వెళ్తాను}}), its \textquotedblleft{}not\textquotedblright{} (\textbf{\te{రాను}} (\emph{rānu}, \textquotedblleft{}I will not come\textquotedblright{})), or a suggestion with \textbf{-\te{దాం}} (\textbf{\te{వెళ్దాం}} (\emph{veḷdāṁ}, \textquotedblleft{}let's go\textquotedblright{})).

\subsection*{You'll want to know: -\te{అకపోతే}, if not}
\begin{quote}
\textbf{\te{నువ్వు} \te{త్వరగా} \te{రాకపోతే}, \te{బస్సు} \te{వెళ్ళిపోతుంది}.}
\end{quote}

\begin{quote}
\emph{nuvvu tvaragā rākapōtē, bassu veḷḷipōtundi.} — \textquotedblleft{}If you do not come quickly, the bus will go.\textquotedblright{}
\end{quote}

\textbf{-\te{అక}} + \textbf{\te{పోతే}}: \textbf{\te{రాకపోతే}}, \textquotedblleft{}if (you) don't come\textquotedblright{}; \textbf{\te{చేయకపోతే}}, \textquotedblleft{}if (you) don't do\textquotedblright{}; \textbf{\te{తినకపోతే}}, \textquotedblleft{}if (you) don't eat\textquotedblright{}. \textbf{\te{లేకపోతే}}, which you know as \textquotedblleft{}otherwise\textquotedblright{}, is the same shape on \textbf{\te{లే}-}, \textquotedblleft{}not be\textquotedblright{}: \textquotedblleft{}if not\textquotedblright{}.

\begin{quote}
\textbf{\te{త్వరగా} \te{రా}, \te{లేకపోతే} \te{బస్సు} \te{వెళ్ళిపోతుంది}.}
\end{quote}

\begin{quote}
\emph{tvaragā rā, lēkapōtē bassu veḷḷipōtundi.} — \textquotedblleft{}Come quickly, otherwise the bus will go.\textquotedblright{}
\end{quote}

\subsection*{Your turn}
Say these aloud:

\begin{itemize}
\raggedright
  \item \textquotedblleft{}If it rains tomorrow, I will not come\textquotedblright{} — \emph{rēpu varṣaṁ paḍitē nēnu rānu}
  \item \textquotedblleft{}If there is time, let's go to a film\textquotedblright{} — \emph{samayaṁ uṇṭē sinimāki veḷdāṁ}
  \item \textquotedblleft{}If you do not come quickly, the bus will go\textquotedblright{} — \emph{nuvvu tvaragā rākapōtē, bassu veḷḷipōtundi}
  \item \textquotedblleft{}Come quickly, otherwise the bus will go\textquotedblright{} — \emph{tvaragā rā, lēkapōtē bassu veḷḷipōtundi}
  \item one plan of yours for tomorrow, with an if and an if not
\end{itemize}

\begin{tcolorbox}[breakable,colback=teal!4,colframe=teal!35!black,title={Before you move on}]
\textquotedblleft{}If it rains\textquotedblright{}? (\textbf{\te{వర్షం} \te{పడితే}}.) \textquotedblleft{}If you don't come\textquotedblright{}? (\textbf{\te{రాకపోతే}}.) \textquotedblleft{}Otherwise\textquotedblright{}? (\textbf{\te{లేకపోతే}}, \textquotedblleft{}if not\textquotedblright{}.) \textquotedblleft{}I will not come\textquotedblright{}? (\textbf{\te{రాను}}.)
\end{tcolorbox}
\section[(review)]{Review — if, and otherwise}

\label{lesson:TE-R324-if-and-otherwise}



\begin{quote}
Say \textquotedblleft{}grief\textquotedblright{}, \textquotedblleft{}a remedy\textquotedblright{}, \textquotedblleft{}one's home place\textquotedblright{}, \textquotedblleft{}a vote\textquotedblright{} and \textquotedblleft{}in the meantime\textquotedblright{}. (\emph{duḥkhaṁ}, \emph{auṣadhaṁ}, \emph{ūru}, \emph{ōṭu}, \emph{īlōgā}.)

\begin{itemize}
\raggedright
  \item \emph{Write it:} \textbf{\te{దుఃఖం}}, \textbf{\te{ఔషధం}} and \textbf{\te{ఊరు}} once each, then \textbf{\te{ఓ}} and \textbf{\te{ఈ}} on their own
\end{itemize}
\end{quote}

\subsection*{Your turn}
Say these aloud:

\begin{itemize}
\raggedright
  \item If you come to my home town, I will be very happy.
  \item If it rains tomorrow, we will not go. Otherwise we will go at six.
  \item If by chance the bus does not come, let's take an auto.
  \item It has been raining for two days. If it rains tomorrow, I will stay at home.
  \item If you read this book, you will like it more than that one.
\end{itemize}

\begin{tcolorbox}[breakable,colback=teal!4,colframe=teal!35!black,title={Before you move on}]
\textquotedblleft{}If\textquotedblright{}? (\textbf{-\te{తే}}: \textbf{\te{వస్తే}}, \textbf{\te{ఉంటే}}, \textbf{\te{పడితే}}.) \textquotedblleft{}If by chance\textquotedblright{}? (\textbf{\te{ఒకవేళ} … -\te{తే}}.) \textquotedblleft{}If not\textquotedblright{}? (\textbf{-\te{అకపోతే}}: \textbf{\te{రాకపోతే}}.) \textquotedblleft{}Otherwise\textquotedblright{}? (\textbf{\te{లేకపోతే}}.)
\end{tcolorbox}
